\documentclass[12pt]{article}

\usepackage[a4paper, total={6in, 8in}]{geometry}
\usepackage{textcomp}

\begin{document}
\setlength{\parindent}{0pt}

\def \t {\textrightarrow}
\def \v {\vspace{1em}}
\def \bi {\begin{itemize}}
\def \ei {\end{itemize}}
\def \s[#1] {\section*{#1}}
\def \ss[#1] {\subsection*{#1}}
\def \sss[#1] {\subsubsection*{#1}}

\s[D'Annunzio - Poesia]
Finora parabola dei romanzi \t ma c'è anche d'annunzio poeta \t scrive le "Laudi" (rimanda a San Francesco d'Assisi) \\
Il sottotitolo è "canzoni di terra ed oltre mare" \t e ci sono 5? \\
Laus = celebrazione, lode \t le lodi sono rivolte a celebrare l'uomo nella natura \\
L'uomo e la natura totalmente umanizzata \t antropoformizzazione della natura e naturalizzazione dell'uomo \t e il mondo è pervaso da un elan vital, slancio vitale = vitalismo d'annunziano \\
Se lo slancio vitale è ovunque, come dio nelle lodi di d'assisi \t si parla di panismo \\
Caratteristica della sua personalita e della sua tempra \t è il suo sensualismo = percezione con i 5 sensi \\
Sono 5: "Maia", "Elettra", "Alcyone", "Asterope" e "Merope" \t Alcione è importante \\
Noi guaridiamo da Alcyone, che esce nel 1903, i testi che ritraggono il volto intimo del poeta \\
La sensualita d'annunziana

\ss[La pioggia nel pineto - Alcyone]
Ermione è una donna \t non è detto che sia la Duse \t ma ci aggancia a una tradizione culturale e di vita che lo porta a una certa lassitudo \t è completamente immerso \\
C'è la declinazione totale delle emozioni dei cinque sensi \t un orchestra sotto innumerevoli dita, che sono le goccie di pioggia, che toccano le piante e gli arbusti e creano una sinfonia di note, che rapiscono i due amanti che si trovano nel pineto

\sss[Fino verso 32]
Si nota un equilibrio dettato dalla ricorrenza di alcuni termini \t che sono costituiti da una sola parola \t si capisce passione per equilibrio tra parola scritta e spazio bianco \\
Questo non ha pero la valenza che ha per ungaretti, e per i futuristi \t qua è una strategia emotiva \t il testo arriva all'animo del lettore anche con questa scansione \\
Prima persona \t l'autore è coinvolto \t l'immersione nella natura ha fatto di loro due creature di vita viventi \\
Enj. con umane \t non / ma odo = antitesi \\
Parlare usato in modo transitivo = suo sperimentalismo, che rapisce l'animo lettore \t e in qualche modo è già rapito da se stesso \\
Le parole sono nuove \t taci e ascolta, specularita \\
La ricorrenza del termine piove = anafora insistente \t le tamerici riporta a pascoli \t sono salmastre e arse perche pineta si trova a ridosso del mare \t a Versiglia probabilmente \\
Il salmastro \t è importante in Montale, per superare d'ann. etc. \t la banderuola salmastra della casa dei doganieri \t anche scagliosi e irti montale \\
Il mirto è molto aromatico \t partecipazione corale di tutti i sensi, e coinvolto anche l'olfatto \\
Fulgente = luminoso \\
Aulesco (latino) = profumare \t effluvi di profumo \\
Silvano = dio delle selve, prima metamorfosi \\
L'anima è novella \t come se tornasse a uno stato primitivo, sulla favola bella che ieri/oggi = chiasmo \t illude con poliptoto \\
Ermione è il nome della figlia di Elena di Troia \t ma lui non partecipa all'ulissismo \t non riprende direttamente Ulisse (anche Foscolo parla di Ulisse)

\sss[Fino verso 64]
Impostato su rimandi \t "Odi" registro è lo stesso, e la pioggia cade in un crepitio (onomatopea) \\
Ascolta \t richiama di nuovo all'attenzione \\
Il pianto australe = pioggia \\
Cè cinerino = cielo scuro, grigio \\
Innumerevoli dita \t le sensazioni \\

\sss[Fino verso 96]
Qua enfasi sull'ascoltare \\
Gli accordi \t parla in termini musicali \\
Le cicale sono aeree \\
Compresenza di sensazioni, suono $\pm$ forte \t da lontano, come un ombra remota \t piu foco... si spegne = climax \\
Nota si spegne = sinestesia \\
Non c'è il suono del mare \t in Boudealire c'è la pioggia battente (in Spleen) \t argentea colore che sembra metallo, che è sinestesia e sineddoche \\
Il suono è diverso a seconda delle fronde \t e le stesse note, suonate diversamente dagli stessi strumenti, si diversificano con la mano dell'artista \t l'emozione fa la differenza
\end{document}
